% 第九卷 9.2.1.4 节：“研究生命”上的词格（教学词典：研究、研究生、生命）。
% 下方为字符LSTM链，上方弧为词典匹配词；词单元从起始字取状态，在词尾字处经门控并入字符单元。
\begin{tikzpicture}[
  x=1mm,y=1mm,figure picture,
  char/.style={rectangle,rounded corners=.8mm,draw=FigureInput,fill=FigureInputFill,
    line width=.6pt,minimum width=10mm,minimum height=7mm,inner sep=0pt,
    font=\figurefont\footnotesize},
  cell/.style={rectangle,rounded corners=.8mm,draw=FigureStroke,fill=white,
    line width=.6pt,minimum width=10mm,minimum height=6mm,inner sep=0pt,
    font=\figurefont\scriptsize},
  word/.style={rectangle,rounded corners=.8mm,draw=FigureModel,fill=white,
    line width=.6pt,minimum height=5.5mm,inner xsep=2pt,font=\figurefont\scriptsize},
  chain/.style={figure flow},
  wordarc/.style={draw=FigureModel,line width=.8pt},
  merge/.style={draw=BookTeal,line width=.9pt,-{Stealth[length=1.6mm,width=1.1mm]}}
]
  \foreach \i/\c/\x in {1/研/0,2/究/22,3/生/44,4/命/66}{
    \node[char] (x\i) at (\x,0) {\c};
    \node[cell] (h\i) at (\x,11) {$\bm c_{\i}$};
    \draw[figure link] (x\i) -- (h\i);
  }
  \draw[chain] (h1) -- (h2);
  \draw[chain] (h2) -- (h3);
  \draw[chain] (h3) -- (h4);
  % 研究：1→2
  \draw[wordarc] (h1.north) to[out=70,in=180] (8,22) -- (14,22);
  \node[word] (wa) at (11,22) {研究};
  \draw[merge] (wa.east) to[out=0,in=110] (h2.north);
  % 研究生：1→3
  \draw[wordarc] ([xshift=-2mm]h1.north) to[out=90,in=180] (14,33);
  \node[word] (wb) at (22,33) {研究生};
  \draw[merge] (wb.east) to[out=0,in=105] ([xshift=-1.5mm]h3.north);
  % 生命：3→4
  \draw[wordarc] ([xshift=1.5mm]h3.north) to[out=75,in=180] (52,24);
  \node[word] (wc) at (55,24) {生命};
  \draw[merge] (wc.east) to[out=0,in=110] (h4.north);
  \node[font=\figurefont\scriptsize,text=FigureMuted,anchor=west] at (74,11) {字符单元};
  \node[font=\figurefont\scriptsize,text=FigureMuted,anchor=west] at (74,0) {输入字符};
  \node[font=\figurefont\scriptsize,text=FigureMuted,anchor=west] at (74,28) {词典词单元};
\end{tikzpicture}
